La promoción de la seguridad operacional en SKYFILMS constituye el conjunto de actividades orientadas a fortalecer la \textbf{cultura de seguridad}, garantizar la \textbf{competencia del personal} y asegurar una \textbf{comunicación efectiva y transversal} de la información en todos los niveles de la organización.\\

El éxito del SMS depende de que cada miembro de la organización comprenda sus responsabilidades, reciba formación adecuada y disponga de canales de comunicación que fomenten la transparencia y el aprendizaje colectivo.

\section{Instrucción y Educación}
SKYFILMS establece un \textbf{programa integral de instrucción y educación en seguridad operacional}, con el objetivo de asegurar que todo el personal involucrado en la operación de UAS cuente con las competencias necesarias para desempeñar sus funciones de manera segura, eficiente y en conformidad con la normativa vigente.\\

\textbf{Principios rectores:}

\begin{itemize}
    \item La instrucción en seguridad operacional será obligatoria para todo el personal que participa directa o indirectamente en las operaciones de UAS.
    \item La formación incluirá \textbf{inducción inicial, entrenamiento recurrente y capacitaciones específicas} asociadas a riesgos identificados en la operación.
    \item El contenido de los programas de capacitación será coherente con los objetivos de seguridad operacional de SKYFILMS, los SPI definidos, y la normativa aplicable (RAC 100 y RAC 219).
    \item Las temáticas de instrucción incluirán de manera obligatoria las \textbf{políticas y procedimientos del explotador UAS}, así como contenidos de \textbf{Factores Humanos (FFHH)} y otras buenas prácticas orientadas a la prevención de errores y a la promoción de la seguridad de la operación.
    \item Se fomentará la participación en \textbf{simulacros de emergencias operacionales y ejercicios prácticos}, con el fin de fortalecer la respuesta del personal ante situaciones críticas.
    \item Los registros de asistencia y evidencias de capacitación se conservarán conforme al sistema de gestión documental del SMS.
\end{itemize}

\textbf{Objetivos específicos de la instrucción:}

\begin{enumerate}
    \item Garantizar el conocimiento y aplicación de los procedimientos de seguridad operacional establecidos.
    \item Promover la \textbf{cultura justa}, asegurando que los empleados comprendan la importancia del reporte voluntario de peligros sin temor a represalias.
    \item Reforzar las competencias técnicas de pilotos, observadores y personal de apoyo en escenarios de riesgo operacional identificados.
    \item Asegurar que la instrucción evolucione de acuerdo con las lecciones aprendidas, los hallazgos de auditorías y las tendencias de seguridad observadas en los SPI.
    \item Asegurar la comprensión de los \textbf{Factores Humanos (FFHH)} y la aplicación de \textbf{buenas prácticas operacionales} como elementos clave en la mejora de la seguridad operacional.
\end{enumerate}

\subsection{Competencias del Personal Clave en Seguridad Operacional}

Con el fin de garantizar la eficacia del Sistema de Gestión de la Seguridad Operacional (SMS), SKYFILMS establecerá y mantendrá definidas las competencias mínimas requeridas para el personal clave en materia de seguridad operacional, asegurando que el programa de instrucción cubra de manera integral las funciones individuales, responsabilidades y obligaciones de rendición de cuentas dentro del sistema.\\

\subsubsection{Competencias del Personal Clave de Seguridad Operacional}

    El personal con responsabilidades directas dentro del SMS deberá demostrar, como mínimo, las siguientes competencias:
    
    \begin{enumerate}
        \item \textbf{Responsable del SMS:}
            \begin{itemize}
                \item Conocimiento sólido de la normativa aplicable (RAC 100, RAC 219 y disposiciones complementarias).

                \item Dominio de los cuatro componentes del SMS.

                \item Capacidad para liderar procesos de identificación de peligros y gestión del riesgo.

                \item Competencia en análisis de datos de seguridad e interpretación de SPI.

                \item Habilidades en auditoría interna y mejora continua.

                \item Capacidad de comunicación efectiva y liderazgo en cultura justa.
            \end{itemize}
        \item \textbf{Personal Operativo (Pilotos UAS, Observadores, Personal de Apoyo):}
            \begin{itemize}
                \item Conocimiento de la Política y Objetivos de Seguridad Operacional.

                \item Comprensión de sus responsabilidades individuales dentro del SMS.

                \item Competencia en procedimientos operacionales normales y de emergencia.

                \item Capacidad para identificar peligros y reportarlos oportunamente.

                \item Aplicación de principios de gestión del riesgo en la toma de decisiones operacionales.
            \end{itemize}
    \end{enumerate}
    
\subsubsection{Cobertura del Programa de Instrucción SMS}

    El programa de instrucción en seguridad operacional garantizará que todo el personal:
        \begin{itemize}
            \item Comprenda cómo funciona el SMS dentro de SKYFILMS.

            \item Conozca los canales formales de reporte de peligros y sucesos.

            \item Identifique claramente sus roles y responsabilidades individuales.

            \item Entienda su obligación de rendición de cuentas frente al cumplimiento de procedimientos.

            \item Reconozca la importancia de la cultura justa y del reporte voluntario.

            \item Comprenda la relación entre objetivos de seguridad, SPI y desempeño operacional.
        \end{itemize}\\

    La instrucción incluirá, como mínimo:

        \begin{itemize}
            \item Estructura y componentes del SMS.

            \item Gestión del riesgo operacional (identificación de peligros, evaluación y mitigación).

            \item Uso del sistema de reporte.

            \item Análisis básico de eventos y lecciones aprendidas.

            \item Responsabilidad individual en la prevención de incidentes.

            \item Factores humanos aplicables a operaciones UAS.

            \item Buenas prácticas operacionales y disciplina operacional.
        \end{itemize}

\textbf{NOTA:} Para la implementación se seguirá el \textbf{Anexo 12 - Programa de Instrucción SMS}.
\subsubsection{Evaluación de Competencias}

    La organización verificará el cumplimiento de las competencias mediante:
        \begin{itemize}
            \item Evaluaciones teóricas en instrucción inicial y recurrente.

            \item Participación en simulacros y ejercicios prácticos.

            \item Revisión de desempeño operacional.

            \item Seguimiento de participación en actividades del SMS.

            \item Retroalimentación derivada de auditorías internas y análisis de SPI.
        \end{itemize}



\section{Comunicación de la Seguridad Operacional}
SKYFILMS fomentará una \textbf{comunicación clara, efectiva y bidireccional} en materia de seguridad operacional, entendida como un elemento esencial para consolidar una cultura organizacional positiva y orientada a la prevención de riesgos.\\

\textbf{Lineamientos de comunicación:}

\begin{itemize}
    \item Establecimiento de un \textbf{sistema de reporte voluntario y confidencial}, accesible para todos los colaboradores, como herramienta clave para la identificación temprana de peligros.
    \item Difusión periódica de \textbf{boletines de seguridad operacional}, que incluyan estadísticas de desempeño, análisis de incidentes, lecciones aprendidas y mejores prácticas.
    \item Realización de \textbf{reuniones de seguridad operacional (Safety Meetings)} lideradas por el Gerente de Seguridad Operacional, con el fin de presentar avances, retroalimentar al personal y reforzar el compromiso de la dirección.
    \item Implementación de \textbf{canales digitales} (intranet, correos institucionales, grupos de trabajo) para garantizar el acceso oportuno a la información de seguridad.
    \item Comunicación inmediata de \textbf{alertas de seguridad operacional} cuando se identifiquen tendencias de riesgo que requieran acciones preventivas urgentes.
    \item Coordinación con terceros, autoridades aeronáuticas y clientes para compartir información relevante que impacte la seguridad de las operaciones.

\end{itemize}

\textbf{NOTA:} En caso de tener información critica a comunicar a las personas involucradas en la operación UAS, el Gerente de Segurar Operacional usara, de los métodos y canales anteriormente mencionados, el método mas eficaz según el contexto.\\

\textbf{Objetivos específicos de la comunicación:}

\begin{enumerate}
    \item Mantener al personal informado sobre los \textbf{resultados del desempeño del SMS} y las acciones derivadas de auditorías, inspecciones y análisis de SPI.
    \item Facilitar la retroalimentación desde todos los niveles de la organización, garantizando que las inquietudes en materia de seguridad sean consideradas en la toma de decisiones.
    \item Generar confianza y compromiso en los colaboradores mediante la transparencia en el manejo de la información.
    \item Promover el aprendizaje organizacional a través de la divulgación de lecciones aprendidas propias y de otros operadores UAS.
\end{enumerate}
